% report_02.tex -- FVCOM-MP late-time spatial diagnostics
% May 2026 | WaterPACT Delaware River Test Cases

\PassOptionsToPackage{table,dvipsnames}{xcolor}
\documentclass[aspectratio=169,10pt]{beamer}
\usetheme{Madrid}
\usecolortheme{beaver}
\setbeamertemplate{navigation symbols}{}
\setbeamerfont{title}{size=\large}

\usepackage{booktabs}
\usepackage{graphicx}
\usepackage{amsmath}
\usepackage{amssymb}
\setlength{\tabcolsep}{4pt}

\newcommand{\ca}{\textcolor{OliveGreen}{\textbf{a}}}
\newcommand{\cb}{\textcolor{NavyBlue}{\textbf{b}}}
\newcommand{\cc}{\textcolor{BurntOrange}{\textbf{c}}}

\title{FVCOM-MP: Late-Time Tracer Diagnostics}
\subtitle{Delaware River Estuary Test Cases}
\author{WaterPACT FVCOM-MP Development}
\date{May 2026}

\begin{document}

\begin{frame}
  \titlepage

  \vspace{-0.8em}
  \begin{center}
    \colorbox{NavyBlue!10}{%
      \parbox{0.78\textwidth}{\centering\small
        Current FVCOM-MP test outputs now span multi-month integrations
        ($>3$ months), enabling day-30 and day-90 tracer checks beyond the
        initial 10-day mass-balance archive.}}
  \end{center}
\end{frame}

\begin{frame}{Delaware River Test Cases}
  \centering
  \vspace{0.3em}
  \footnotesize
  \begin{tabular}{lcccc}
    \toprule
    & & Case \ca & Case \cb & Case \cc \\
    \cmidrule(lr){3-5}
    Role & & \textit{baseline (published)} & sensitivity b & sensitivity c \\
    \midrule
    \texttt{PLAST\_FLOC}   & & F & T & T \\
    \texttt{SEDIMENT\_MODEL} & & F$^{\dagger}$ & T & T \\
    \texttt{SED\_SD50}     & & --- & 0.1\,mm & 1.0\,mm \\
    MP water field(s)      & & $c$ &
                             \multicolumn{2}{c}{$c_a + c_d$} \\
    Agg.\ gate $d_{MP}^{\max}$ & & --- &
                             $34\,\mu$m (excl.) &
                             $162\,\mu$m (act.) \\
    MP settling est. & & --- &
                             \shortstack{$c_d\approx-0.39$\,mm\,s$^{-1}$\\
                             buoyant/free MP} &
                             \shortstack{$c_a=+3.0$\,mm\,s$^{-1}$\\
                             $\Delta\approx3.4$\,mm\,s$^{-1}$} \\
    \bottomrule
  \end{tabular}

  \vspace{0.5em}
  {\scriptsize
    $c$\,=\,disaggregated conc.\hfill
    $c_a$\,=\,aggregated,\;$c_d$\,=\,disaggregated (floc path)\\[2pt]
    $^{\dagger}$sed.\ initialized; \texttt{adv\_scal}/FCT bypassed
    $\;\bullet\;$
    \texttt{SOURCE\_AG\_RATE\,=\,0} $\Rightarrow$ $c_a\equiv 0$
    $\;\bullet\;$
    $\rho_p=950\,\mathrm{kg\,m^{-3}}$,\; $d_p=100\,\mu\mathrm{m}$
    $\;\bullet\;$
    settling estimate uses $\rho_w\approx1025\,\mathrm{kg\,m^{-3}}$ and
    \texttt{SED\_WSET}=3\,mm\,s$^{-1}$}
\end{frame}

\begin{frame}{First 10-Day Mass Balance: v012 Ramp-Corrected Analytic Archive}
  \centering
  \includegraphics[width=0.92\textwidth,height=0.72\textheight,keepaspectratio]{%
    ../FVCOM_MP_test_run/PYTHON/output/mass_conservation/archive/%
    v012_ramp_corrected_analytic/total_mass_by_case.png}

  \vspace{0.2em}
  \colorbox{OliveGreen!14}{%
    \parbox{0.82\textwidth}{\centering\small
      Final 10-day archive check: total mass by case is well behaved after the
      ramp-corrected analytic source treatment.}}
\end{frame}

\begin{frame}{Day 30 Spatial Contrast: Water Column and Bed Mass}
  \centering
  \scriptsize
  \begin{tabular}{ccc}
    \textbf{Case a: \texttt{mp1\_davg}} &
    \textbf{Case b: \texttt{mp1\_davg}} &
    \textbf{Case c: \texttt{mp1\_davg}} \\
    \includegraphics[width=0.30\textwidth,height=0.275\textheight,keepaspectratio]{%
      ../FVCOM_MP_test_run/PYTHON/output/spatial_distribution/day30/planview_dayavg_mp1_davg_casea.png} &
    \includegraphics[width=0.30\textwidth,height=0.275\textheight,keepaspectratio]{%
      ../FVCOM_MP_test_run/PYTHON/output/spatial_distribution/day30/planview_dayavg_mp1_davg_caseb.png} &
    \includegraphics[width=0.30\textwidth,height=0.275\textheight,keepaspectratio]{%
      ../FVCOM_MP_test_run/PYTHON/output/spatial_distribution/day30/planview_dayavg_mp1_davg_casec.png} \\
    \textbf{Case a: bed mass} &
    \textbf{Case b: bed mass} &
    \textbf{Case c: bed mass} \\
    \includegraphics[width=0.30\textwidth,height=0.275\textheight,keepaspectratio]{%
      ../FVCOM_MP_test_run/PYTHON/output/spatial_distribution/day30/planview_dayavg_bot_mass_casea.png} &
    \includegraphics[width=0.30\textwidth,height=0.275\textheight,keepaspectratio]{%
      ../FVCOM_MP_test_run/PYTHON/output/spatial_distribution/day30/planview_dayavg_bot_mass_caseb.png} &
    \includegraphics[width=0.30\textwidth,height=0.275\textheight,keepaspectratio]{%
      ../FVCOM_MP_test_run/PYTHON/output/spatial_distribution/day30/planview_dayavg_bot_mass_casec.png}
  \end{tabular}
\end{frame}

\begin{frame}{Floc-Component Split: Case c, Day 30 vs. Day 90}
  \centering
  \scriptsize
  \begin{tabular}{cc}
    \textbf{Day 30: \texttt{mp1\_agg\_davg}} &
    \textbf{Day 30: \texttt{mp1\_dis\_davg}} \\
    \includegraphics[width=0.39\textwidth,height=0.305\textheight,keepaspectratio]{%
      ../FVCOM_MP_test_run/PYTHON/output/spatial_distribution/day30/planview_dayavg_mp1_agg_davg_casec.png} &
    \includegraphics[width=0.39\textwidth,height=0.305\textheight,keepaspectratio]{%
      ../FVCOM_MP_test_run/PYTHON/output/spatial_distribution/day30/planview_dayavg_mp1_dis_davg_casec.png} \\
    \textbf{Day 90: \texttt{mp1\_agg\_davg}} &
    \textbf{Day 90: \texttt{mp1\_dis\_davg}} \\
    \includegraphics[width=0.39\textwidth,height=0.305\textheight,keepaspectratio]{%
      ../FVCOM_MP_test_run/PYTHON/output/spatial_distribution/day90/planview_dayavg_mp1_agg_davg_casec.png} &
    \includegraphics[width=0.39\textwidth,height=0.305\textheight,keepaspectratio]{%
      ../FVCOM_MP_test_run/PYTHON/output/spatial_distribution/day90/planview_dayavg_mp1_dis_davg_casec.png}
  \end{tabular}

  \vspace{0.2em}
  {\small Case \cc\ shows the active floc pathway; case \cb\ remains effectively
  zero in the aggregated component under the fine-sediment parameterization.}
\end{frame}

\end{document}
